\section{Limitations}\label{sec:limits}
\begin{itemize}
\item \textbf{The floor is cross-document variability, not instrument error.} It includes systematic differences between assays and laboratories and curation errors. Relative to a single well-controlled assay it overstates the irreducible error; relative to a model that can learn document-specific offsets (random splits) it may also be too pessimistic, as the cases with $f>1$ show. We have not yet evaluated document-grouped splits, which would separate these effects.
\item \textbf{Dependence and representativeness of replicates.} The floor is estimated from the 1.8--11.1\% of molecules that were measured in several documents. These are more potent on average than other molecules in several tasks, are concentrated in few document pairs for some tasks (one pair contributes 79\% of molecule pairs for CHEMBL237 $\mathrm{IC}_{50}$), and include identical values that are probably re-publications. We assume one noise variance per task. Shared authors, re-published values and assay conditions (substrate or ATP concentration, species, mutants) were not checked.
\item \textbf{Sensitivity to curation choices.} The estimated noise depends on the trimming threshold, on how duplicated values and stereoisomers are treated and on assay selection (Table~\ref{tab:robust}); for some tasks these choices change $\hat s$ by 0.1 or more. The assay confidence threshold of 8 admits homologous proteins.
\item \textbf{Leakage.} Decontamination covers MoleculeNet and TDC but not Polaris datasets or the pretraining corpora of the foundation models, which are large subsets of PubChem and ZINC and may contain test molecules. Such overlap would be expected to favour the pretrained models, which nevertheless perform worse. Removing MoleculeNet and TDC molecules changed test errors very little (Section~\ref{sec:results}).
\item \textbf{Model coverage.} We test two small chemical language models, frozen, on all tasks and only a D-MPNN and one fine-tuned model on five tasks with three seeds; larger or more recent foundation models, other heads and longer hyperparameter searches may behave differently. The hyperparameter budget (three configurations) is small for all models.
\item \textbf{Splits.} The temporal split leaks some information (ties at the cut-off year; training labels include later documents), and the inner validation split used for hyperparameters is random for all outer splits.
\item \textbf{Scope.} Potency endpoints ($K_i$, $\mathrm{IC}_{50}$) from ChEMBL only; we make no claim about ADMET or physicochemical endpoints. $K_i$ and $\mathrm{IC}_{50}$ tasks also differ in target class (mostly GPCRs versus mostly kinases), so differences between them are not attributable to the measurement type alone. Test labels are noisy; we do not correct test metrics for attenuation.
\item \textbf{Literature.} Our account of related work is based on the published abstracts and metadata of several cited papers; we did not reproduce their analyses.
\end{itemize}
